Some performance plots obtained from the \starlight\ HIJING Overlay MC
  sample are shown in Figure~\ref{fig:pbpb_mc_performance}, for muons 
  passing the Medium working point (Section~\ref{sec:muon_selection}),
  and in Figure~\ref{fig:pbpb_mc_performance_tight}, for muons 
  passing the Tight working point.
In general it is observed that the $\eta$ and $\phi$ positions and the 
  \pt\ values of the reconstructed and true muon are very well correlated.


\begin{figure}[h]
\centering
\includegraphics[width=0.45\textwidth]%
{figures/AllFigs/can_ProbMatch_medium}%
\includegraphics[width=0.45\textwidth]%
{figures/AllFigs/can_pTMatch_medium}
\includegraphics[width=0.45\textwidth]%
{figures/AllFigs/can_EtaMatch_medium}
\includegraphics[width=0.45\textwidth]%
{figures/AllFigs/can_PhiMatch_medium}
\caption{
Some of the MC performance plots.
The top-left panel shows the generated-match propability for
  all reconstructed muons, passing the Medium working point.
The remaining panels show the difference in \pt\ $\eta$ and $\phi$ between
  generated and reconstructed muons for muons with a match-probability>0.5.
Plots are made for MC sample with HIJING overlay.
\label{fig:pbpb_mc_performance}
}
\end{figure}

\begin{figure}[h]
\centering
\includegraphics[width=0.45\textwidth]%
{figures/AllFigs/can_ProbMatch_tight}%
\includegraphics[width=0.45\textwidth]%
{figures/AllFigs/can_pTMatch_tight}
\includegraphics[width=0.45\textwidth]%
{figures/AllFigs/can_EtaMatch_tight}
\includegraphics[width=0.45\textwidth]%
{figures/AllFigs/can_PhiMatch_tight}
\caption{
Same as Figure~\ref{fig:pbpb_mc_performance} but for Tight muons.
\label{fig:pbpb_mc_performance_tight}
}
\end{figure}
\FloatBarrier


